This chapter shows the exact specifications for the \glspl{crd} implemented, how they are validated and how their fields affect the resources.
We shall go over every \gls{crd} individually, show the specification and detail the important fields.
The fields shown here will be limited to the spec part of the resource.
Fields in the status block shall, if necessary, be explained in the following chapter.
In the light of promise theory, every resource here can be thought of as a promise to the user, specifying behavior of the operator.
Kubernetes in this case acts as the platform, forwarding the promise of the operator made with its \gls{crd} to the user of the cluster.

\subsection{API Version in Kubernetes}

In order to implement an \gls{api} in Kubernetes, there needs to be separation between resources, versioning and naming of the resources.
Kubernetes groups objects into \enquote{API versions}, which consist of an identifier for the \gls{api} as well as a version itself.
In this case, the \enquote{ansible-operator.lightjack.de/v1alpha1} \gls{api} version was created.
Here, the \enquote{v1alpha1} indicates the first version of the operator \enquote{v1} with an alpha state resource \enquote{alpha} in the first version \enquote{1}.
Should there be breaking changes to the \gls{api} once the operator is in a stable condition, this version should be changed and the resources moved to a different \gls{api} version.

\subsection{The \anhost resource}
The complete \anhost type declaration can be found in the appendix under section \ref{app:anhostdeclare}.

A \anhost may be defined as follows:

%TC:ignore
\begin{lstlisting}[language=yaml,caption=Example \anhost definition]
apiVersion: ansible-operator.lightjack.de/v1alpha1
kind: AnsibleHost
metadata:
  name: ssh-node-1
spec:
  ansibleName: ssh-node-1
  connection:
    host: ssh-node-1.default.svc.cluster.local
    user: ansible
  ssh:
    sshKeySecretRef:
      name: ansible-operator-hosts-private-key
    sshHostKeySecretRef:
      name: ssh-node-1-host-key
    ignoreHostKey: false
  privilege:
    become: true
\end{lstlisting}
%TC:endignore

One notable variable is the \enquote{AnsibleName} which is used to specify the name of the resource within ansible.

